\documentclass{ltxdoc}
\usepackage{simplex}

\usepackage{microtype}

\newcommand{\env}[1]{\texttt{#1}}

\title{The \textsf{simplex} package\\[4pt]
       \large Linguistic example environments for \LaTeX}
\author{Kyle Gorman\\\small\texttt{kylebgorman@gmail.com}}
\date{Version 1.1}

\EnableCrossrefs
\CodelineIndex
\RecordChanges

\begin{document}
\maketitle

\begin{abstract}
The \textsf{simplex} package provides three numbered example environments
for use in linguistics papers---\env{unlabeledexample}, \env{shortexample},
and \env{example}---together with interlinear gloss macros (\cs{gll},
\cs{glll}, \cs{glt}, \cs{gln}, \cs{glend}) adapted from
\textsf{covington.sty} and \textsf{gloss.tex}.
\end{abstract}

\tableofcontents

\section{Introduction}

Linguistics papers routinely include numbered examples, sometimes with
descriptive labels, and frequently with word-by-word interlinear glosses
paired with a free translation. The \textsf{simplex} package provides a
minimal, self-contained set of macros for both tasks.

The gloss alignment code is adapted from \textsf{gloss.tex}
(v1.0, Marcel R.~van der Goot, 1991) and \textsf{covington.sty}
(Michael Covington). The example counter environments are original.

\section{Requirements}

\textsf{simplex} requires \LaTeXe{} and automatically loads the
\textsf{xparse} package (part of the standard \LaTeX{} distribution).

\section{Package options}
\label{sec:options}

\DescribeMacro{resetchapter}
The \texttt{resetchapter} option resets the \texttt{exampleno} counter
at the start of each \cs{chapter}, so that numbering restarts at
(1) in every chapter. It requires a document class that
defines \cs{chapter} (e.g.\ \textsf{report} or \textsf{book}).

\begin{verbatim}
\usepackage[resetchapter]{simplex}
\end{verbatim}

Without this option (the default), \texttt{exampleno} simply
increments throughout the document, as in v1.0. Enabling it also
changes what \cs{exref} (\S\ref{sec:examples}) prints for
cross-chapter references; see below.

\section{Numbered example environments}
\label{sec:examples}

All three environments share a single counter, \texttt{exampleno}, which
increments automatically and is typeset as an Arabic numeral in
parentheses, e.g.\ (1), (2), \ldots

\DescribeEnv{unlabeledexample}
The \env{unlabeledexample} environment produces a numbered example with
no label. It takes no arguments.

\begin{verbatim}
\begin{unlabeledexample}
The cat sat on the mat.
\end{unlabeledexample}
\end{verbatim}

\begin{unlabeledexample}
The cat sat on the mat.
\end{unlabeledexample}

\DescribeEnv{shortexample}
The \env{shortexample} environment produces a numbered example whose
optional label appears on the \emph{same line} as the number. The label
is given as an optional argument in square brackets.

\begin{verbatim}
\begin{shortexample}[\textsc{Substance-freeness of structural change}]
Features changed by rule application need not be present in the rule's
structural environment.
\end{shortexample}
\end{verbatim}

\begin{shortexample}[\textsc{Substance-freeness of Structural Change}]
Features changed by rule application need not be present in the rule's
structural environment.
\end{shortexample}

\DescribeEnv{example}
The \env{example} environment is like \env{shortexample} except that the
label appears on a \emph{separate line} above the example content,
separated by a small vertical skip. This is useful when the label is long
or when the content is multi-line (such as a gloss block).

\begin{verbatim}
\begin{example}[Assibilating perfect passive participles and agents]
\begin{tabular}[H]{llllll}
metere & `reap' & messus & `harvested' & messor & `reaper' \\
fodere & `dig'  & fossus & `dug'       & fossor & `digger' \\
\end{tabular}
\end{example}
\end{verbatim}

\begin{example}[Assibilating perfect passive participles and agents]
\begin{tabular}[H]{llllll}
metere & `reap' & messus & `harvested' & messor & `reaper' \\
fodere & `dig'  & fossus & `dug'       & fossor & `digger' \\
\end{tabular}
\end{example}

\subsection{Cross-referencing}

All three environments call \cs{refstepcounter}, so standard
\cs{label}/\cs{ref} cross-referencing works normally.

\begin{verbatim}
\begin{unlabeledexample}
The cat sat on the mat.
\label{ex:basic}
\end{unlabeledexample}

See (\ref{ex:basic}).
\end{verbatim}

\DescribeMacro{\exref}
\cs{exref} is an alternative to \cs{ref} purpose-built for citing
examples: it adds the surrounding parentheses itself, and---when the
\texttt{resetchapter} option (\S\ref{sec:options}) is in
effect---it automatically includes the chapter number for references
that cross a chapter boundary.

\begin{verbatim}
See \exref{ex:basic}.
\end{verbatim}

With \texttt{resetchapter} \emph{not} given, \verb|\exref{ex:basic}|
simply prints (1), matching \verb|(\ref{ex:basic})|. With
\texttt{resetchapter} given, a reference to an example in the
\emph{current} chapter still prints as (1), but a reference
to an example in chapter~2 cited from a later chapter prints as
(2/1).

\cs{exref} also takes an optional argument that postposes a letter
inside the parentheses, for referring to a sub-part of an example:

\begin{verbatim}
See \exref[a]{ex:basic}.
\end{verbatim}

\noindent which prints (1a), or (2/1a) for a
cross-chapter reference under \texttt{resetchapter}.

\cs{exref} reads the label's stored value directly rather than via
\cs{ref}, so it is unaffected by packages (such as \textsf{hyperref})
that redefine \cs{ref}; it does not currently add a hyperlink of its
own.

\section{Interlinear gloss macros}
\label{sec:gloss}

\textsf{simplex} provides macros for typesetting word-aligned interlinear
glosses of the kind common in linguistic fieldwork and typological work.
The macros accept two or three lines of space-separated tokens and stack
them vertically, aligning each column.

\subsection{Two- and three-line glosses: \texttt{\textbackslash gll} and \texttt{\textbackslash glll}}

\DescribeMacro{\gll}
\cs{gll} introduces a two-line gloss: a source-language line followed by
a morpheme-gloss line. The two lines are separated by a literal newline
(end of line) inside the macro argument; each line ends with \verb|\\|.
After the two gloss lines, supply a translation with \cs{glt} (or
\cs{gln} if no translation is needed), and optionally close with
\cs{glend}.

\DescribeMacro{\glll}
\cs{glll} is like \cs{gll} but accepts three aligned lines (typically:
source, morpheme gloss, and a second-tier gloss or transliteration).

\begin{verbatim}
\begin{example}[\textsc{Pre-Liquid Shortening} bled by external sandhi]
\gll clāmōr=ad cael-um                
     shout=to  heaven-\textsc{acc.sg.}
\glt `a shout to heaven' %(fragments of Ennius)
% Omitted---no-op here.
%\glend
\end{example}
\end{verbatim}

\begin{example}[\textsc{Pre-Liquid Shortening} bled by external sandhi]
\gll clāmōr=ad cael-um                
     shout=to  heaven-\textsc{acc.sg.}
\glt `a shout to heaven' %(fragments of Ennius)
% Omitted---no-op here.
%\glend
\end{example}

\subsection{Translation and closing macros}

\DescribeMacro{\glt}
\cs{glt} closes the aligned gloss block and introduces a free
translation. It must appear after the final gloss line even when no
translation is needed (pass an empty line in that case).

\DescribeMacro{\gln}
\cs{gln} is an alternative to \cs{glt} for use when no translation is
wanted at all; it closes the gloss block without adding a newline.

\DescribeMacro{\glend}
\cs{glend} is a semantic no-op provided for readability; it marks the
end of a gloss block.

\subsection{Customizing gloss fonts}

Three token-list hooks control the font applied to each line:

\begin{center}
\begin{tabular}{ll}
\cs{eachwordone} & font for the source line (default: \cs{rm})\\
\cs{eachwordtwo} & font for the first gloss line (default: \cs{rm})\\
\cs{eachwordthree} & font for the second gloss line (default: \cs{rm})\\
\end{tabular}
\end{center}

To italicize the source line, for example, add to your preamble:

\begin{verbatim}
\let\eachwordone=\it
\end{verbatim}

The inter-word spacing in gloss blocks is controlled by the skip register
\cs{glossglue}, whose default value is \texttt{5pt plus 2pt minus 1pt}.

\section{Known limitations}

\begin{itemize}
  \item Per-chapter numbering requires the \texttt{resetchapter}
        option (\S\ref{sec:options}); it is off by default for
        backward compatibility with v1.0 documents.
  \item \cs{gll} and \cs{glll} rely on active end-of-line characters and
        may behave unexpectedly inside certain environments or
        macro arguments.
  \item Long words without spaces will not be broken across lines.
\end{itemize}

\section{License}

This package is copyright \copyright\ Kyle Gorman and is distributed
under the \LaTeX{} Project Public License, version~1.3 or later
(\url{http://www.latex-project.org/lppl.txt}).
The gloss alignment code in \texttt{simplex.sty} was adapted from 
\textsf{gloss.tex} by Marcel R.~van der Goot (1991).
Additional adaptations are derived from \textsf{covington.sty} by
Michael Covington.

\section{Version history}

\begin{description}
  \item[v1.1 (2026/06/30)] Added the \texttt{resetchapter} package
    option and the \cs{exref} cross-referencing macro (with optional
    postposed letter).
  \item[v1.0 (2026/03/02)] Initial public release.
\end{description}

\PrintIndex

\end{document}
